\section{Dataset}

The primary dataset for this project is HAM10000, a public benchmark of dermoscopic pigmented skin lesion images released for academic machine-learning research through the ISIC Archive and Harvard Dataverse~\cite{tschandl2018ham10000}. The standard benchmark contains 10,015 RGB dermoscopic images across seven diagnostic categories: actinic keratoses and intraepithelial carcinoma (\texttt{akiec}), basal cell carcinoma (\texttt{bcc}), benign keratosis-like lesions (\texttt{bkl}), dermatofibroma (\texttt{df}), melanocytic nevi (\texttt{nv}), melanoma (\texttt{mel}), and vascular lesions (\texttt{vasc}).

HAM10000 was selected because it is public, multi-class, feasible in Google Colab, and difficult enough to make the required feature-extraction, fusion, and fine-tuning comparisons meaningful. The dataset is class-imbalanced, so this project treats macro-F1 and per-class F1 as important analysis metrics in addition to the assignment-required accuracy, precision, recall, and F1-score.

To reduce leakage risk, the intended split policy is stratified train/validation/test splitting with lesion-level grouping when lesion identifiers are available. This avoids placing near-duplicate images of the same lesion in both training and test partitions.

EuroSAT is the backup dataset if the medical-image workflow becomes impractical. It provides 27,000 Sentinel-2 image patches across 10 land-use classes and is operationally simpler while still satisfying the public multi-class image classification requirement~\cite{helber2019eurosat}.
